% Slides do acompanhamento — sem nenhuma métrica de teste, conforme o enunciado.
\documentclass[aspectratio=169,10pt,t]{beamer}
\newcommand{\tipodoc}{Acompanhamento}
% Tema comum dos slides (acompanhamento e pitch).
\usepackage{fontspec}
\setsansfont{Noto Sans}
\setmonofont[Scale=0.85]{Hack}
\usepackage[none]{hyphenat}
\usepackage{booktabs,tabularx,multirow,array,graphicx,xcolor}
\usepackage{tikz}
\usepackage{newunicodechar}
% Noto Sans não tem estes símbolos; usa a versão matemática.
\newunicodechar{→}{\ensuremath{\rightarrow}}
\newunicodechar{≈}{\ensuremath{\approx}}
\newunicodechar{≥}{\ensuremath{\geq}}
\newunicodechar{×}{\ensuremath{\times}}
\newunicodechar{↔}{\ensuremath{\leftrightarrow}}
\newcolumntype{L}{>{\raggedright\arraybackslash}X}
% Macros comuns ao artigo e aos slides.
\definecolor{azul}{HTML}{1B3A5C}
\definecolor{azulclaro}{HTML}{2E6B8F}
\definecolor{cinza}{HTML}{5F6B76}
\definecolor{pendcor}{HTML}{B3261E}
\definecolor{pendfundo}{HTML}{FBEDEC}

% Valores gerados por scripts/dados_docs.py; o que ainda não existe aparece como campo em branco.
\newcommand{\resdef}[2]{\expandafter\def\csname res@#1\endcsname{#2}}
\newcommand{\vazio}{\textcolor{pendcor}{\rule[-1.5pt]{2.4em}{0.5pt}}}
\newcommand{\res}[1]{\ifcsname res@#1\endcsname\csname res@#1\endcsname\else\vazio\fi}
\InputIfFileExists{gerado/dados.tex}{}{}

% Trecho a completar depois dos treinos (aparece em vermelho para não passar despercebido).
\newcommand{\pend}[1]{\textcolor{pendcor}{[#1]}}

% \figuraarq{largura}{altura}{arquivo em resultados/figuras}{o que gera a figura}
\newcommand{\figdir}{../../resultados/figuras}
\newcommand{\figuraarq}[4]{%
  \IfFileExists{\figdir/#3}%
    {\includegraphics[width=#1,height=#2,keepaspectratio]{\figdir/#3}}%
    {\fcolorbox{pendcor}{pendfundo}{\parbox[c][#2][c]{\dimexpr#1-2\fboxsep-2\fboxrule}{\centering\small
      \textcolor{pendcor}{\textbf{Figura pendente}}\\[2pt]\texttt{resultados/figuras/}\\\texttt{\detokenize{#3}}\\[2pt]#4}}}}


\renewcommand{\figurename}{Figura}
\renewcommand{\tablename}{Tabela}

\setbeamersize{text margin left=9mm,text margin right=9mm}
\setbeamertemplate{navigation symbols}{}
\setbeamercolor{normal text}{fg=black!85}
\setbeamercolor{structure}{fg=azul}
\setbeamercolor{alerted text}{fg=azulclaro}
\setbeamerfont{frametitle}{size=\Large,series=\bfseries}
\setbeamertemplate{itemize item}{\color{azulclaro}\raisebox{1.5pt}{\rule{4pt}{4pt}}}
\setbeamertemplate{itemize subitem}{\color{cinza}--}
\setbeamertemplate{enumerate item}{\color{azulclaro}\bfseries\insertenumlabel.}
\setlength{\leftmargini}{1.3em}

% Quem apresenta cada slide (participação equilibrada).
\newcommand{\falante}{}
\newcommand{\fala}[1]{\renewcommand{\falante}{#1}}

\setbeamertemplate{frametitle}{%
  \vskip6mm
  \begin{beamercolorbox}[wd=\paperwidth,leftskip=9mm,rightskip=9mm]{frametitle}
    \usebeamerfont{frametitle}\color{azul}\insertframetitle
    \ifx\falante\empty\else\hfill{\normalsize\mdseries\color{cinza}\falante}\fi
  \end{beamercolorbox}
  \vskip-1mm{\color{azulclaro}\hspace*{9mm}\rule{18mm}{1.6pt}}\vskip2mm}

\setbeamertemplate{footline}{%
  \hbox to\paperwidth{\hspace*{9mm}\color{cinza}\tiny Projeto 3 \textperiodcentered{} Classificação de Cenas Aéreas \textperiodcentered{} \tipodoc\hfill\insertframenumber\hspace*{9mm}}\vskip3mm}

\setbeamertemplate{title page}{%
  \vfill
  {\color{azulclaro}\small\bfseries CIÊNCIA DE DADOS \textperiodcentered{} UNIDADE III \textperiodcentered{} UNIFOR\par}
  \vskip3mm
  {\color{azul}\fontsize{26}{30}\selectfont\bfseries\inserttitle\par}
  \vskip2mm
  {\color{azulclaro}\large\insertsubtitle\par}
  \vskip3mm{\color{azulclaro}\rule{25mm}{2pt}}\vskip5mm
  {\normalsize\insertauthor\par}
  \vskip2mm
  {\small\color{cinza}\insertdate\par}
  \vfill}

\newcommand{\repo}{\href{https://github.com/JoeyAlanS/redes-reurais-fotos-satelite}{github.com/JoeyAlanS/redes-reurais-fotos-satelite}}
% Caixa de destaque
\newcommand{\destaque}[1]{\colorbox{azul!7}{\parbox{\dimexpr\linewidth-2\fboxsep}{\raggedright\small #1}}}


\title{Classificação de Cenas Aéreas com Redes Neurais}
\subtitle{Acompanhamento \textperiodcentered{} dados, preparação, arquiteturas e primeiros resultados de validação}
\author{Paulo Henrico Rabelo (2312652) \quad Joey Alan (2320416)}
\date{Projeto 3 \textperiodcentered{} NWPU-RESISC45 \textperiodcentered{} \repo}

\begin{document}

\begin{frame}[plain]
\titlepage
\end{frame}

\fala{Paulo}
\begin{frame}{Problema e perguntas}
\begin{columns}[T]
\begin{column}{0.55\textwidth}
\textbf{Desafio:} como a arquitetura e a resolução afetam a classificação de cenas aéreas, e quais ganhos justificam o custo?
\vskip4mm
\begin{enumerate}
\item MLP A/B × CNN A/B em 128\,×\,128
\item Efeito de 256\,×\,256 em cada CNN
\item O ganho compensa o custo?
\item Quais classes se confundem, e por quê?
\end{enumerate}
\end{column}
\begin{column}{0.4\textwidth}
\destaque{\textbf{6 treinamentos}, todos com pesos novos:\\[2pt]
MLP-A-128 \quad MLP-B-128\\
CNN-A-128 \quad CNN-B-128\\
CNN-A-256 \quad CNN-B-256}
\vskip3mm
\destaque{Seleção pelo \textbf{F1 macro de validação}. O teste só é avaliado no fim, uma única vez.}
\end{column}
\end{columns}
\end{frame}

\begin{frame}{Dados filtrados: 10 classes do NWPU-RESISC45}
\begin{columns}[T]
\begin{column}{0.42\textwidth}
\begin{itemize}
\item \res{total} imagens RGB 256\,×\,256, 700 por classe
\item Cópia pública no Hugging Face (o OneDrive deu HTTP 403), conferida por SHA-256
\item Todas abertas e verificadas: \textbf{\res{problemas}} problemas de leitura, dimensão ou canal
\item Ordem fixa das classes em código, métricas e matrizes
\item Originais nunca alterados
\end{itemize}
\end{column}
\begin{column}{0.55\textwidth}
\centering\figuraarq{\linewidth}{5.2cm}{eda_exemplos.png}{notebooks/01\_exploracao.ipynb}
\end{column}
\end{columns}
\end{frame}

\begin{frame}{Duplicatas e prevenção de vazamento}
\begin{columns}[T]
\begin{column}{0.44\textwidth}
\begin{itemize}
\item \textbf{Exatas} (SHA-256): \res{dup/exatas} pares, todos em \textit{airport}
\item \textbf{Quase-duplicatas}: dHash de 256 bits nas 8 orientações, limiar de \res{dup/limiar} bits
\begin{itemize}
\item \textit{mountain} 111/324: mesma imagem recodificada
\item \textit{dense\_residential} 175/638: cena deslocada
\end{itemize}
\item Demais imagens: vizinha mais próxima a \textbf{≥ 35 bits}
\item Cada grupo é sorteado como \textbf{uma unidade}
\end{itemize}
\end{column}
\begin{column}{0.53\textwidth}
\centering\figuraarq{\linewidth}{5.2cm}{eda_duplicatas.png}{notebooks/01\_exploracao.ipynb}
\end{column}
\end{columns}
\end{frame}

\begin{frame}{Partições: 70/15/15, estratificadas, semente 42}
\begin{columns}[T]
\begin{column}{0.4\textwidth}
\centering
\begin{tabular}{@{}lrrr@{}}
\toprule
& Treino & Valid. & Teste\\
\midrule
Por classe & 490 & 105 & 105\\
\textit{airport} & \res{part/airport/treino} & \res{part/airport/validacao} & \res{part/airport/teste}\\
\midrule
\textbf{Total} & \textbf{\res{part/total/treino}} & \textbf{\res{part/total/validacao}} & \textbf{\res{part/total/teste}}\\
\bottomrule
\end{tabular}
\vskip4mm
\begin{itemize}\small
\item Caminhos ordenados antes do sorteio
\item \texttt{particoes.csv} reutilizado em tudo, inclusive no bônus
\item Exploração visual só no treino
\end{itemize}
\end{column}
\begin{column}{0.56\textwidth}
\centering\figuraarq{\linewidth}{5.4cm}{eda_contagens.png}{notebooks/01\_exploracao.ipynb}
\end{column}
\end{columns}
\end{frame}

\fala{Joey}
\begin{frame}{Exploração: onde cor e textura não bastam}
\begin{columns}[T]
\begin{column}{0.44\textwidth}
Centróide mais próximo em cor e textura, só no treino:
\begin{itemize}
\item Floresta \textbf{89\%} e deserto \textbf{76\%}: bem separados
\item Montanha \textbf{14\%}: vai para rio (27\%), agrícola circular (19\%) e aeroporto (14\%)
\item Aeroporto se confunde com rio (22\%)
\end{itemize}
\vskip2mm
\destaque{Hipótese: essas classes exigem \textbf{forma e organização espacial}, o que favorece a CNN.}
\end{column}
\begin{column}{0.53\textwidth}
\centering\figuraarq{0.9\linewidth}{5.2cm}{eda_confusao_centroides.png}{notebooks/01\_exploracao.ipynb}
\end{column}
\end{columns}
\end{frame}

\begin{frame}{Mesma imagem em 256\,×\,256 e 128\,×\,128}
\begin{columns}[T]
\begin{column}{0.36\textwidth}
\begin{itemize}
\item Redimensionada \textbf{a partir do original} (bilinear com \textit{antialias}), pixels ÷ 255
\item Em 128 px cada pixel junta 2\,×\,2 do original
\item Tanques, casas e barcos perdem contorno
\item 256 px custa \textbf{4×} mais pixels por camada convolucional
\end{itemize}
\end{column}
\begin{column}{0.6\textwidth}
\centering\figuraarq{\linewidth}{5.2cm}{eda_resolucao.png}{notebooks/01\_exploracao.ipynb}
\end{column}
\end{columns}
\end{frame}

\begin{frame}{Arquiteturas e hipóteses}
\small
\begin{tabularx}{\linewidth}{@{}lLrL@{}}
\toprule
Modelo & Estrutura & Parâmetros & Hipótese\\
\midrule
MLP A & 49.152 → Dense 256 → 10 & 12.586.506 & referência sem viés espacial\\
MLP B & 49.152 → 256 → 128 → 64 → 10 & 12.626.314 & mesma ordem de parâmetros: isola a \textbf{profundidade}\\
CNN A & 3 blocos Conv-BN-ReLU-Pool (32/64/128) → GAP → 10 & 95.210 & padrões locais com 132× menos parâmetros\\
CNN B & 4 blocos (32/64/128/256) → GAP → Dense 128 → 10 & 424.042 & campo receptivo de 22 → 46 px\\
\bottomrule
\end{tabularx}
\vskip4mm
\begin{itemize}
\item GAP: os parâmetros das CNNs são \textbf{idênticos} em 128 e 256 px, e nenhum peso é reaproveitado entre resoluções
\item Em 256 px, o campo receptivo da CNN A cobre só 8,6\% do lado da cena; espera-se que a B ganhe mais
\end{itemize}
\end{frame}

\begin{frame}{Protocolo de treinamento}
\begin{columns}[T]
\begin{column}{0.5\textwidth}
\begin{itemize}
\item TensorFlow/Keras, pesos novos, semente 42
\item Entropia cruzada, Adam (0,001), lote 32
\item Até 40 épocas; para após \textbf{8 sem melhora}
\item Checkpoint = época de maior \textbf{F1 macro de validação}
\item Aumento só no treino: espelhamentos e giros de 90°
\end{itemize}
\end{column}
\begin{column}{0.46\textwidth}
\destaque{\textbf{Isolar a resolução:} em cada CNN, mesmo pré-processamento, semente, otimizador, taxa, lote e parada nas duas resoluções.}
\vskip3mm
\destaque{Cada treino registra tempo total e por época, tamanho, hardware e versões das bibliotecas (\texttt{modelos/*.json}).}
\end{column}
\end{columns}
\end{frame}

\fala{Paulo}
\begin{frame}{Problema encontrado e corrigido}
\begin{columns}[T]
\begin{column}{0.52\textwidth}
Num teste de 2 épocas com o \textit{BatchNorm} padrão do Keras:
\begin{itemize}
\item Treino: \textbf{66\%} de acurácia; validação: \textbf{11\%}, com quase tudo previsto como uma classe só
\item Variância móvel da 1ª camada: 0,048; real nos lotes: 0,003
\item Causa: com \textit{momentum} 0,99, a média móvel ainda guardava o valor inicial 1 (0,99\textsuperscript{306} ≈ 0,05)
\end{itemize}
\end{column}
\begin{column}{0.44\textwidth}
\destaque{\textbf{Correção:} \textit{momentum} 0,9. As estatísticas de inferência convergem na 1ª época.}
\vskip3mm
\destaque{Mesmo teste após a correção: validação de 41\% na 1ª época e 67\% na 2ª (só para verificar, não é resultado final).}
\end{column}
\end{columns}
\end{frame}

\fala{Paulo e Joey}
\begin{frame}{Resultados de validação (teste não avaliado)}
\begin{columns}[T]
\begin{column}{0.42\textwidth}
\centering\small
\begin{tabular}{@{}lrrr@{}}
\toprule
Modelo & Acur. & F1 macro & Época\\
\midrule
MLP-A-128 & \res{MLP-A-128/acuracia_validacao} & \res{MLP-A-128/f1_macro_validacao} & \res{MLP-A-128/checkpoint_epoca}\\
MLP-B-128 & \res{MLP-B-128/acuracia_validacao} & \res{MLP-B-128/f1_macro_validacao} & \res{MLP-B-128/checkpoint_epoca}\\
CNN-A-128 & \res{CNN-A-128/acuracia_validacao} & \res{CNN-A-128/f1_macro_validacao} & \res{CNN-A-128/checkpoint_epoca}\\
CNN-B-128 & \res{CNN-B-128/acuracia_validacao} & \res{CNN-B-128/f1_macro_validacao} & \res{CNN-B-128/checkpoint_epoca}\\
CNN-A-256 & \res{CNN-A-256/acuracia_validacao} & \res{CNN-A-256/f1_macro_validacao} & \res{CNN-A-256/checkpoint_epoca}\\
CNN-B-256 & \res{CNN-B-256/acuracia_validacao} & \res{CNN-B-256/f1_macro_validacao} & \res{CNN-B-256/checkpoint_epoca}\\
\bottomrule
\end{tabular}
\vskip3mm
\raggedright\footnotesize Mínimo exigido: uma MLP e uma CNN. Tempo por época: \pend{preencher}.
\end{column}
\begin{column}{0.55\textwidth}
\centering\figuraarq{\linewidth}{5.2cm}{painel_curvas_perda.png}{scripts/prever.py e scripts/avaliar.py (só validação)}
\end{column}
\end{columns}
\end{frame}

\fala{Joey}
\begin{frame}{Próximos passos}
\begin{enumerate}
\item Concluir os seis treinamentos no mesmo hardware (GPU) e salvar checkpoints e históricos
\item Comparar pela validação e fechar as escolhas, sem olhar o teste
\item Avaliar os seis modelos salvos no teste, uma única vez: métricas por classe, matrizes, custo e erros
\item Discutir as confusões à luz da exploração (montanha, rio, aeroporto)
\item Escrever o artigo (até 10 páginas, ABNT) e preparar o pitch
\item Bônus opcional: YOLO \textit{classify} pré-treinado nas mesmas partições
\end{enumerate}
\end{frame}

\fala{Paulo}
\begin{frame}{Dúvidas}
\begin{enumerate}
\item O critério de quase-duplicatas (dHash nas 8 orientações, limiar de 12 bits) é suficiente, ou esperam outra verificação?
\item Usar o \textbf{mesmo aumento de dados} (espelhos e giros de 90°) também nas MLPs é adequado para a comparação?
\item Selecionar e parar pelo F1 macro de validação, em vez da perda, está de acordo com o protocolo?
\item Treinar no Google Colab em sessões diferentes, mas na mesma GPU (T4), conta como "mesmo hardware"?
\item No artigo, o espaçamento simples é aceito para caber em 10 páginas?
\end{enumerate}
\vfill
\centering\small Código e partições: \repo
\end{frame}

\end{document}
